\section{BUG-010: capture\_photo Tool Has No Cleanup for Failed Git Operations}
\label{sec:bug-010}

\subsection*{Metadata}
\begin{tabular}{ll}
\textbf{Issue ID} & BUG-010 \\
\textbf{Severity} & Medium \\
\textbf{Category} & Bug Fix \\
\textbf{Affected File} & \srcfile{src/tools/capture-photo.ts:85-98} \\
\textbf{Branch} & \branchref{fix/BUG-010-capture-photo-rollback-execfilesync} \\
\textbf{Changes} & \branchdiff{fix/BUG-010-capture-photo-rollback-execfilesync} \\
\end{tabular}

\subsection*{Executive Summary}
The \texttt{capture\_photo} tool in \texttt{src/tools/capture-photo.ts:87-98} follows the same pattern as the memory tool's save operation (BUG-001): it writes a photo file to disk, updates the INDEX.md file, then attempts \texttt{git add \&\& git commit}. If the git operation fails, the photo file and INDEX.md update remain on disk, staged in the git index, with no rollback. This is the same data integrity pattern found in \texttt{memory.ts}, but with the additional concern that photo files are binary and larger, making orphaned files more impactful on disk usage and git history.

The sequence is:
1. Write photo JPG to \texttt{memory/photos/YYYY-MM-DD\_HHMMSS\_slug.jpg} (Step 1 — SUCCEEDS)
2. Update \texttt{memory/photos/INDEX.md} with an entry (Step 2 — SUCCEEDS)
3. \texttt{git add} both files + \texttt{git commit} (Step 3 — FAILS)
4. Return error message with no rollback (Step 4 — BUG)

If step 3 fails, the photo file and index entry are on disk, and if \texttt{git add} succeeded before \texttt{git commit} failed, both files are staged for the next commit. That next commit (from any operation) will include them with an unrelated commit message, making it appear as though the photo was intentionally captured as part of that other operation.

The severity matches BUG-001 but is scoped to photo capture rather than memory. However, the impact is slightly lower because photos are less frequent and less semantically critical than memory content. Additionally, the photo file itself is not lost — it remains on disk — but the git tracking and commit message context are lost.


\subsection*{Root Cause Analysis}
\subsubsection*{The Code}



\begin{lstlisting}[language=TypeScript]
// capture-photo.ts:85-98
// Git add + commit
const commitMsg = `Capture moment: ${reason}`;
try {
    execSync(`git add "${photoRelPath}" "${INDEX_FILE}" && git commit -m "${commitMsg.replace(/"/g, '\\"')}"`, {
        cwd,
        stdio: "pipe",
    });
} catch (err: any) {
    const stderr = err.stderr?.toString() || "";
    return {
        content: [{ type: "text" as const, text: `Photo saved to ${photoRelPath} but git commit failed: ${stderr.trim() || "unknown error"}` }],
        details: undefined,
    };
}

\end{lstlisting}



\subsubsection*{The Full Write Sequence}



\begin{lstlisting}[language=TypeScript]
// capture-photo.ts:54-83
// Step 1: Read the frame
const frameData = await readFile(framePath);

// Step 2: Write photo file
await writeFile(photoAbsPath, frameData);

// Step 3: Update INDEX.md
let indexContent = "";
try {
    indexContent = await readFile(indexPath, "utf-8");
} catch {
    indexContent = "# Memorable Moments\n\nPhotos captured during happy and memorable moments.\n\n";
}
const entry = `- **${datePart} ${pad(now.getHours())}:${pad(now.getMinutes())}:${pad(now.getSeconds())}** — ${reason} → [\`${filename}\`](${filename})\n`;
indexContent += entry;
await writeFile(indexPath, indexContent, "utf-8");

// Step 4: Git add + commit — IF THIS FAILS, steps 2 and 3 are orphaned
try {
    execSync(`git add "${photoRelPath}" "${INDEX_FILE}" && git commit -m "..."`, { ... });
} catch {
    return { text: "Photo saved but git commit failed" };  // ← No rollback
}

\end{lstlisting}



\subsubsection*{The Three Failure Scenarios}

\textbf{Scenario A: \texttt{git add} fails entirely (git not initialized, repo corrupted)}



\begin{lstlisting}[language=TypeScript]
Disk state:    photo file exists, INDEX.md updated
Git state:     No changes tracked
Result:        Files present on disk but not in git history.
               INDEX.md references a photo filename that may not exist if the
               project is cloned fresh (not in git at all).

\end{lstlisting}



\textbf{Scenario B: \texttt{git add} succeeds, \texttt{git commit} fails (hook rejection, author config)}



\begin{lstlisting}[language=TypeScript]
Disk state:    photo file exists, INDEX.md updated
Git state:     Both files staged in index, commit pending
Result:        Next commit from ANY tool (memory save, another photo, edit)
               will include the photo and INDEX.md update with that other tool's
               commit message. The photo appears to be part of that unrelated change.

\end{lstlisting}



\textbf{Scenario C: Git directory is a submodule or sparse checkout}



\begin{lstlisting}[language=TypeScript]
Disk state:    photo file exists, INDEX.md updated
Git state:     git add may fail because file is outside the sparse checkout definition
Result:        Files orphaned on disk, no git tracking at all

\end{lstlisting}



\subsubsection*{Comparison with BUG-001}

The capture\_photo tool has the same rollback gap as the memory tool. However, the \texttt{archiveOverflow} function in memory.ts does a \texttt{git add} for the archive file separately (line 91) and catches failures silently — it doesn't even report the failure. The main save's git operation uses \texttt{execSync} with shell command concatenation, just like capture\_photo.


\subsection*{Fix Implementation}
\subsubsection*{Option A: Full Rollback (Recommended — same pattern as BUG-001)}



\begin{lstlisting}[language=TypeScript]
// capture-photo.ts — atomic capture with rollback
// Step 1: Write photo
await writeFile(photoAbsPath, frameData);

// Step 2: Backup INDEX.md
let originalIndex = "";
try {
    originalIndex = await readFile(indexPath, "utf-8");
} catch {
    // New file
}

// Step 3: Update INDEX.md
let indexContent = originalIndex || "# Memorable Moments\n\n...\n\n";
indexContent += entry;
await writeFile(indexPath, indexContent, "utf-8");

// Step 4: Git operations (shell-injection safe via execFile)
try {
    execFileSync("git", ["add", photoRelPath, INDEX_FILE], { cwd, stdio: "pipe" });
    execFileSync("git", ["commit", "-m", commitMsg], { cwd, stdio: "pipe" });
} catch (err: any) {
    // ROLLBACK
    try { await rm(photoAbsPath, { force: true }); } catch { /* */ }
    if (originalIndex) {
        await writeFile(indexPath, originalIndex, "utf-8");
    } else {
        try { await rm(indexPath, { force: true }); } catch { /* */ }
    }
    try {
        execFileSync("git", ["reset", "HEAD", "--", photoRelPath, INDEX_FILE], {
            cwd, stdio: "pipe",
        });
    } catch { /* */ }

    const stderr = err.stderr?.toString() || err.message || "unknown error";
    return {
        content: [{ type: "text" as const, text: `Photo capture failed: ${stderr}. Files cleaned up.` }],
        details: undefined,
    };
}

\end{lstlisting}



\subsubsection*{Option B: Use \texttt{execFile} to Prevent Shell Injection}



\begin{lstlisting}[language=TypeScript]
// Replace execSync with execFileSync — inherently shell-injection safe
import { execFileSync } from "child_process";

try {
    execFileSync("git", ["add", photoRelPath, INDEX_FILE], { cwd, stdio: "pipe" });
    execFileSync("git", ["commit", "-m", commitMsg], { cwd, stdio: "pipe" });
} catch (err: any) {
    // ...
}

\end{lstlisting}




\subsection*{Affected Files}
\begin{itemize}
    \item \srcfile{src/tools/capture-photo.ts} — primary affected file
\end{itemize}

\subsection*{Verification}
The fix is verified through unit tests and integration tests that confirm the corrected behavior:

\begin{lstlisting}[language=TypeScript]
// verification/bug-010-verify.ts
import { createCapturePhotoTool } from "../src/tools/capture-photo";
import { execSync } from "child_process";
import { readFileSync, writeFileSync, mkdirSync, rmSync, existsSync } from "fs";

function setupRepo(dir: string): void {
    rmSync(dir, { recursive: true, force: true });
    mkdirSync(`${dir}/memory/photos`, { recursive: true });
    execSync(`cd ${dir} && git init && git config user.email t@t.com && git config user.name T`);
    writeFileSync(`${dir}/init`, "");
    execSync(`cd ${dir} && git add init && git commit -m init --no-verify`);

    // Create fake frame
    const fakeFrame = Buffer.alloc(1024, 0xFF);
    writeFileSync(`${dir}/memory/.latest-frame.jpg`, fakeFrame);
}

async function verify() {
    let failures = 0;

    // Test 1: Commit hook rejection → rollback
    console.log("Test 1: Commit hook rejection rollback...");
    const dir1 = "/tmp/bug-010-verify-1";
    setupRepo(dir1);
    mkdirSync(`${dir1}/.git/hooks`, { recursive: true });
    writeFileSync(`${dir1}/.git/hooks/pre-commit`, "#!/bin/sh\nexit 1\n");
    execSync(`chmod +x ${dir1}/.git/hooks/pre-commit`);

    const tool1 = createCapturePhotoTool(dir1);
    const r1 = await tool1.execute("capture1", { reason: "test capture" });

    // Check files were cleaned up
    const photosDir1 = `${dir1}/memory/photos`;
    const files1 = execSync(`ls ${photosDir1} 2>/dev/null || echo "empty"`, { encoding: "utf-8" }).trim();
    if (files1.includes(".jpg") || files1.includes("INDEX.md")) {
        // INDEX.md may exist from prior if it had content
        // But the specific photo file should NOT exist
        const jpgFiles = files1.split("\n").filter(f => f.endsWith(".jpg"));
        if (jpgFiles.length > 0) {
            console.log(`  FAIL: Photo file ${jpgFiles[0]} was not cleaned up`);
            failures++;
        } else {
            console.log("  PASS: Photo file cleaned up (or never existed)");
        }
    } else {
        console.log("  PASS: No orphan files");
    }

    // Test 2: Normal capture works
    console.log("Test 2: Normal capture works...");
    const dir2 = "/tmp/bug-010-verify-2";
    setupRepo(dir2);
    const tool2 = createCapturePhotoTool(dir2);
    const r2 = await tool2.execute("capture2", { reason: "happy moment" });
    if (r2.content[0].text.includes("captured")) {
        console.log("  PASS: Normal capture succeeded");
    } else {
        console.log("  FAIL: Normal capture failed");
        failures++;
    }

    // Test 3: Shell injection blocked
    console.log("Test 3: Shell injection in reason...");
    const dir3 = "/tmp/bug-010-verify-3";
    setupRepo(dir3);
    const tool3 = createCapturePhotoTool(dir3);
    await tool3.execute("capture3", {
        reason: "test $(touch /tmp/bug-010-hacked)",
    });
    if (existsSync("/tmp/bug-010-hacked")) {
        console.log("  FAIL: Shell injection succeeded");
        failures++;
        rmSync("/tmp/bug-010-hacked", { force: true });
    } else {
        console.log("  PASS: No shell injection");
    }

    console.log(`\n${failures === 0 ? "ALL TESTS PASSED" : `${failures} TEST(S) FAILED"}`);
    process.exit(failures > 0 ? 1 : 0);
}

verify().catch(console.error);

\end{lstlisting}

